Microservice architecture has been the latest trend in building cloud-native applications and more and more companies have chosen to migrate from the so-called monolithic architecture to microservice architecture~\cite{ICSA17MSTrend, Zhou_TSE18MS, Guo_FSE20Industry}.
Industrial microservice systems can include hundreds to thousands of services.
For example, the e-commerce system of Alibaba contains more than 30,000 services and includes over 300 subsystems.
Microservice systems support the core business of the companies, especially for Internet companies which provide online services through the Internet.

Availability issues have been a critical challenge for large-scale industrial microservice systems.
These systems are highly dynamic and complex.
A service can have several to thousands of instances running on different containers and dynamically created or
destroyed according to the scaling requirements at runtime;
the execution of a microservice system may involve a huge number of microservice interactions and most of them are asynchronous and involve complex invocation chains~\cite{Zhou_TSE18MS}.
In these systems, any anomaly with the quality (e.g., performance, reliability) of a service may propagate along service call chains and finally cause availability issues at the business level (e.g., drop of successfully placed orders).
When an availability issue is detected, its root cause service and anomaly type need to be located in a short time (e.g., 3 minutes) to allow the developers to fix the issue quickly.

Existing approaches cannot efficiently support the root cause localization of availability issues for large-scale microservice systems.
Developers in industry often rely on visualization tools to analyze logs and traces to identify possible anomalies and anomaly propagation chains to locate root causes~\cite{Zhou_TSE18MS, Guo_FSE20Industry}.
Researchers have proposed approaches for automatic root cause localization for microservice or the broader service-based systems using trace analysis~\cite{Cuong_traceBased, FSE2019LatentError} or service dependency graph~\cite{cscs2013monitorRank, CauseInfer2014, Hanzhang_GraphBased2019, GraphBasedformicroservice, NOMS2020MicroRCA, DCCGRID2018CloudRange, ICSOC2018Microscope}.
Trace analysis based approaches require expensive collection and processing of trace data, thus cannot efficiently work for large-scale systems.
Service dependency graph based approaches construct service dependency graphs based on service calls and causal relationships (e.g., services co-located in the same machines).
These approaches locate root causes by traversing service dependency graphs and detecting possible anomalies with the quality metrics (e.g., response time) of services.
These approaches are limited due to their inaccurate detection of service anomalies and inefficient traversing of service dependency graphs, especially when the system has many services and dependencies.

In this paper, we propose a high-efficient root cause localization approach for availability issues of microservice systems, called \app.
Given an availability issue initially reported on a service (called \textit{initial anomalous service}), \app~locates the root cause services and anomaly types (e.g., performance anomaly, reliability anomaly, and traffic anomaly) that cause the issue.
It dynamically constructs a service call graph of the target microservice system, which reflects the service dependencies and related quality metrics (e.g., response time) in the latest time window.
Based on the graph, it analyzes possible anomaly propagation chains from the initial anomalous service by traversing the graph along anomalous service calls.
To detect possible anomalies with individual service calls, we design an anomaly detection model for each of the three popular anomaly types (i.e., performance, reliability, traffic).
To improve the efficiency of anomaly propagation chain analysis, \app~adopts a pruning strategy to eliminate anomalous service call edges that are irrelevant to the current anomaly propagation chain.
Finally, \app~ranks candidate root causes based on the correlations between their quality metrics and the business metrics of the initial anomalous service.

To evaluate the effectiveness and efficiency of \app , we conduct a series of experimental studies with the data and availability issues collected from the e-commerce system of Alibaba.
The results show that \app~significantly outperforms two state-of-the-art baseline approaches MonitorRank~\cite{cscs2013monitorRank} and Microscope~\cite{ICSOC2018Microscope} in terms of both the accuracy and efficiency.
The results also confirm the effectiveness of the pruning strategy, which can significantly improve the efficiency while keeping the accuracy.

\app~has been deployed in Alibaba for more than 5 months and used to handle more than 600 availability issues.
It achieves a top-3 hit ratio of 68\% and reduces the time of typical root cause localization from 30 minutes to 5 minutes.
Feedback from the developers shows that most of them trust the recommendations of \app~and treat the recommendations as the root causes by default.
The analysis of the results also suggests some improvements of the approach.

The remainder of this paper is organized as follows.
Section II introduces background knowledge about the problem scenarios.
Section III presents an overview of \app~and Section IV details the anomaly propagation chain analysis.
Section V describes our experimental study for the evaluation of \app.
Section VI introduces the practical application of \app~in Alibaba.
Section VII discusses related work and Section VIII concludes the paper.


%Related work is discussed in Section II. Section III gives an overview of the AlarmCause system. Root cause localization procedures are detailed in Section IV. Section V describes the experimental evaluation and Section VI concludes the paper.
%The rest of this paper is structured as follows. Section~\ref{sec:background} introduces the background of our work. Section~\ref{sec:overview} presents the overview of \app~, and the detail of \app~is given in Section~\ref{sec:propagation}. Section~\ref{sec:experimental} analysis the experimental result of \app. The practical application is demonstrated in Section~\ref{sec:practical}. Section~\ref{sec:related} reviews the related work, and Section~\ref{sec:conclusion} summarizes the conclusions and describes planned work.
